\documentclass[10pt]{article}

\usepackage[utf8]{inputenc}
\usepackage[margin=0.68in]{geometry}
\usepackage{graphicx}
\usepackage{booktabs}
\usepackage{amsmath}
\usepackage{amssymb}
\usepackage{array}
\usepackage{caption}

\graphicspath{{./}{supplementary_assets_v2/}{supplementary_assets/}}
\captionsetup{font=footnotesize,labelfont=bf}
\setlength{\textfloatsep}{7pt}
\setlength{\floatsep}{7pt}
\setlength{\intextsep}{7pt}
\setlength{\abovecaptionskip}{3pt}
\setlength{\belowcaptionskip}{2pt}

\newcommand{\pp}{percentage points}

\begin{document}

\begin{center}
{\Large \textbf{Supplementary Material}}\\[0.35em]
{\large \textbf{MAML-Select: An Online Adaptive Client Selection Method for Federated Learning via Meta-Learning}}\\[0.25em]
{\normalsize \textit{Concise revision supplement}}
\end{center}

\section*{S1. Experimental Protocol}

This supplement records the additional evidence used to support the revised letter.  The focus is deliberately compact: reproducibility settings, MAML-Select resource savings, statistical support, sensitivity and ablation evidence, and selector-scaling behaviour.  The complete baseline table is retained in the main manuscript; here we highlight the results most directly connected to the reviewers' requests.

\begin{table}[!ht]
\centering
\caption{Shared experimental configuration.}
\label{tab:supp_config}
\small
\begin{tabular}{@{}ll@{}}
\toprule
Component & Setting \\
\midrule
Clients / selected clients & $N=100$, $K=10$ \\
Communication rounds & $200$ for Fashion-MNIST/CIFAR-10; $150$ for CIFAR-100 analysis \\
Local training & $E=5$ epochs, batch size $32$, client LR $0.01$ \\
Data heterogeneity & Dirichlet non-IID split, $\alpha=0.5$ \\
Models & LightCNN-9 for Fashion-MNIST; ResNet-18 for CIFAR-10/100 \\
MAML-Select policy & MLP $6$--$64$--$64$--$1$ with $4{,}673$ parameters \\
Meta optimization & One inner look-ahead step, $\beta=0.01$; Adam outer update, $\eta=0.001$ \\
Default trade-off weight & $\lambda=0.5$ \\
Seeds & $42$, $123$, $2026$ \\
\bottomrule
\end{tabular}
\end{table}

The heterogeneous-client model uses three device tiers with fixed seeded assignment: low-end clients ($20\%$), mid-range clients ($50\%$), and high-end clients ($30\%$).  Per-round wall-clock time is straggler-bound by the slowest selected client.  Energy is accumulated from the selected clients' tier power and modeled training time, and carbon is obtained from the declared grid intensity of $475$~g\,CO$_2$/kWh.  Since all methods select $K=10$ clients per round, MAML-Select's resource improvements come from adaptive client choice rather than from reducing the number of participating clients.

\section*{S2. Resource and Statistical Evidence}

Table~\ref{tab:supp_reductions} summarizes MAML-Select relative to FedAvg on the shared seeds.  Across all three datasets, MAML-Select keeps full participation coverage while reducing cumulative training compute, energy, and carbon.  The savings are largest on Fashion-MNIST and CIFAR-10, and remain measurable on CIFAR-100 despite the harder model and evaluation setting.

\begin{table}[!ht]
\centering
\caption{MAML-Select resource profile relative to FedAvg over shared seeds.  Positive reductions indicate lower resource use for MAML-Select.}
\label{tab:supp_reductions}
\small
\begin{tabular}{@{}lrrrrr@{}}
\toprule
Dataset & Acc. (\%) & TFLOPs $\downarrow$ & Energy $\downarrow$ & Carbon $\downarrow$ & Coverage \\
\midrule
Fashion-MNIST & $90.11 \pm 0.47$ & $12.8\%$ & $16.4\%$ & $16.4\%$ & $100\%$ \\
CIFAR-10 & $75.63 \pm 1.43$ & $21.5\%$ & $24.7\%$ & $24.7\%$ & $100\%$ \\
CIFAR-100 & $58.15 \pm 0.51$ & $1.7\%$ & $4.7\%$ & $4.7\%$ & $100\%$ \\
\bottomrule
\end{tabular}
\end{table}

\begin{table}[!ht]
\centering
\caption{Paired statistical support for MAML-Select resource reductions versus FedAvg.  Reported $p$ values are two-sided paired $t$-tests over shared seeds; $|d_z|$ is the standardized paired effect size.}
\label{tab:supp_stats}
\small
\begin{tabular}{@{}llrrr@{}}
\toprule
Dataset & Metric & Mean reduction & $|d_z|$ & $p$ \\
\midrule
Fashion-MNIST & TFLOPs & $17.93$ & $5.93$ & $0.009$ \\
 & Energy (Wh) & $943.58$ & $6.74$ & $0.007$ \\
 & Carbon (g) & $448.20$ & $6.74$ & $0.007$ \\
\midrule
CIFAR-10 & TFLOPs & $1792.16$ & $12.68$ & $0.002$ \\
 & Energy (Wh) & $1187.44$ & $18.11$ & $0.001$ \\
 & Carbon (g) & $564.04$ & $18.11$ & $0.001$ \\
\midrule
CIFAR-100 & TFLOPs & $107.80$ & $5.15$ & $0.012$ \\
 & Energy (Wh) & $169.65$ & $10.89$ & $0.003$ \\
 & Carbon (g) & $80.58$ & $10.89$ & $0.003$ \\
\bottomrule
\end{tabular}
\end{table}

\begin{figure}[!ht]
\centering
\includegraphics[width=0.92\textwidth]{supp_resource_trajectories}
\caption{Round-wise cumulative compute, energy, and carbon trajectories.  The MAML-Select curve remains resource-efficient throughout training, showing that the reported savings are accumulated over the whole FL process rather than appearing only at the final round.}
\label{fig:supp_resource}
\end{figure}

\section*{S3. Sensitivity and State-Feature Ablation}

The sensitivity study evaluates the cost trade-off weight $\lambda\in\{0.1,0.5,1.0,5.0\}$ on Fashion-MNIST.  The selected value $\lambda=0.5$ gives a balanced operating point: accuracy remains close to the best observed value while compute, energy, and carbon are reduced relative to the low-penalty setting.  The ablation study verifies that the six-dimensional state vector gives stable behaviour; removing one feature at a time changes accuracy only modestly, which indicates that the learned selector is not dependent on a single fragile signal.

\begin{figure}[!ht]
\centering
\begin{minipage}{0.49\textwidth}
\centering
\includegraphics[width=\linewidth]{fig_lambda_sensitivity_multimetric_boxed}
\caption*{(a) $\lambda$ sensitivity.}
\end{minipage}
\hfill
\begin{minipage}{0.49\textwidth}
\centering
\includegraphics[width=\linewidth]{fig_feature_ablation_trimmed_boxed}
\caption*{(b) State-feature ablation.}
\end{minipage}
\caption{Design evidence for MAML-Select.  The $\lambda$ sweep shows the accuracy--resource trade-off, and the ablation plot shows that the policy remains stable under single-feature removals.}
\label{fig:supp_lambda_ablation}
\end{figure}

\begin{table}[!ht]
\centering
\caption{Fashion-MNIST sensitivity and ablation summaries.}
\label{tab:supp_design}
\footnotesize
\begin{minipage}{0.48\textwidth}
\centering
\begin{tabular}{@{}lrrr@{}}
\toprule
$\lambda$ & Acc. (\%) & TFLOPs & Energy (Wh) \\
\midrule
0.1 & $90.40 \pm 0.22$ & 128.3 & 5176 \\
0.5 & $90.23 \pm 0.55$ & 121.6 & 4799 \\
1.0 & $89.74 \pm 0.79$ & 118.6 & 4639 \\
5.0 & $89.37 \pm 1.02$ & 107.8 & 4140 \\
\bottomrule
\end{tabular}
\end{minipage}
\hfill
\begin{minipage}{0.48\textwidth}
\centering
\begin{tabular}{@{}lrr@{}}
\toprule
Variant & Acc. (\%) & TFLOPs \\
\midrule
Full state & $90.23 \pm 0.55$ & 121.6 \\
w/o loss & $90.16 \pm 0.66$ & 122.8 \\
w/o grad. norm & $90.30 \pm 0.24$ & 122.4 \\
w/o latency & $90.14 \pm 0.14$ & 118.9 \\
w/o battery & $89.90 \pm 0.44$ & 122.2 \\
w/o frequency & $90.47 \pm 0.13$ & 115.0 \\
w/o staleness & $90.37 \pm 0.19$ & 122.7 \\
\bottomrule
\end{tabular}
\end{minipage}
\end{table}

\section*{S4. Selector Scaling}

The selector-scaling probe runs MAML-Select for $50$ rounds while increasing $N=\{20,40,80,100\}$ and keeping $K/N=0.1$.  The measured per-round selection overhead stays in the millisecond range and is small compared with local training time.  This supports the complexity argument in the main paper: scoring all available clients with a fixed-size MLP is linear in the client pool, $O(Nd)$, and selecting the top-$K$ clients adds only the ranking step.

\begin{figure}[!ht]
\centering
\includegraphics[width=0.63\textwidth]{fig_scaling_maml_only_boxed}
\caption{MAML-Select client-selection overhead as the pool size increases.}
\label{fig:supp_scaling}
\end{figure}

\begin{table}[!ht]
\centering
\caption{Mean MAML-Select selection overhead in seconds.}
\label{tab:supp_scaling}
\small
\begin{tabular}{@{}lrrrr@{}}
\toprule
Method & $N=20$ & $N=40$ & $N=80$ & $N=100$ \\
\midrule
MAML-Select & 0.012689 & 0.011160 & 0.011907 & 0.011871 \\
\bottomrule
\end{tabular}
\end{table}

\section*{S5. Summary}

The supplementary evidence strengthens the revised manuscript in four ways.  First, the experimental settings and resource model are fully specified.  Second, MAML-Select shows consistent reductions in modeled compute, energy, and carbon while retaining full client coverage.  Third, the resource reductions are supported by paired statistical tests over the shared seeds.  Fourth, the sensitivity, ablation, and scaling studies show that the learned selector is tunable, stable under state-feature perturbations, and inexpensive to execute.

\end{document}
